\subsection{Smooth sources after imposing a third kernel value}
\label{sec:DE-A4-contact-third-value}

\begin{lemma}[The moving discriminant spans the contact coefficient space]
\label{lem:DE-A4-contact-discriminant-span}
\status{Proved; root and monodromy full independent proof/code audits passed.}
Keep the homogeneous contact space $\Lambda=\mathbf P^4_{[F:T:U:v:h]}$
of \eqref{eq:ss-contact-linear-system}, for any two nonzero points
$g_1=(a,b)$ and $g_2=(c,d)$ of $E:y^2+y=x^3$.
The second point may be a flex.
There is an irreducible hypersurface
$\Delta_{\rm mov}\subset\Lambda$ such that the locus of singular
source curves is contained in
\[
              (FThv=0)\ \cup\ \Delta_{\rm mov}.
\]
The hypersurface $\Delta_{\rm mov}$ is not contained in any
projective hyperplane.  In particular, every hyperplane of
$\Lambda$ distinct from $F=0,T=0,h=0,v=0$ has a nonempty open
parametrizing globally smooth source curves.
\end{lemma}
\begin{proof}
We first check that the source-singularity support calculation
extends to the flex case; none of its multiplicity calculations
will be needed.  Put
\[
 X=x+a,\quad Y=y+b,\quad Z=z+c,\quad
 \ell=w+c^2z+d^2,\quad A_1=Y+Xx^2.
\]
The five functions of the contact system are
$X,Y,\ell,Xz,X\ell$.  Its base points are
$P=(O,O)$, $Q=(g_1,g_2)$ and $(g_1,-2g_2)$, with the last two
coinciding when $g_2$ is a flex.
On the affine open $X\ne0$, the ratios
$Y/X,\ell/X,z,\ell$ detect both tangent directions except where
$A_1=\ell=0$.  This is a finite set: $A_1$ is a nonconstant
function of the first elliptic point and $\ell$ of the second.
At those exceptions the value and the derivative of $z$ still
give first-jet rank at least two.

On the fibre through $g_1$ the section is $T\ell$.
Away from its base points, a singularity therefore lies in $T=0$.
At $Q$ the only first derivative is $v$.  If the additional base
point $(g_1,-2g_2)$ is distinct, its two first derivatives impose
independent conditions; if it coincides with $Q$, it is already
accounted for by $v=0$.  On the first-factor fibre through $-g_1$
the value and first derivatives have rank three except at finitely
many second points, and rank at least two at those exceptions.
On the first-factor fibre through $O$, the value is $F$, except
at $P$, where the first derivative is $h$.  On the second-factor
fibre through $O$, the value and first derivative are $T+hX,h$;
the full first-jet rank is three except for the two finite points
with $X=0$, where it is at least two.
These are exactly the four-chart calculations in the proof of
Lemma~\ref{lem:ss-uniform-source-smoothness}; they remain valid
when the two affine base points merge.
Every finite rank-two exceptional set contributes a singularity
incidence of dimension at most two, and hence no additional divisor
in the four-dimensional coefficient space.

The remaining singularity incidence is a $\mathbf P^1$-bundle over
the irreducible rank-three jet open of the surface.  Its image
closure is irreducible.  It is a hypersurface because the following
explicit singular sections have an ordinary node at the selected
moving point.  On the dense open $XA_1\ell Z\ne0$, solve the
three singularity equations with free parameters $T,h$:
\[
 \begin{aligned}
 F&=\frac{T\ell}{A_1},\\
 U&=Z^2\left(\frac TX+h\right),\\
 v&=\frac{TZ^3}{X}
       +h\bigl(z^2Z+w+d\bigr)+FX^2.
 \end{aligned}
\]
These are the earlier moving-node formulas, after replacing
$E_0$ by $U+hc^2$.  The mixed Hessian is $TZ^2/X$ and is nonzero
when $T\ne0$.  The two critical-point equations then determine
the critical point formally, and its value is a nonzero first-order
linear condition on the coefficients.  Thus the image has codimension
one.  Call it $\Delta_{\rm mov}$.  The remaining incidence strata
have smaller images, and their closures are already contained in
the source discriminant.  The same dimension bounds show that a
generic member of $\Lambda$ is smooth.  Thus the full-system
discriminant hypersurface of
Proposition~\ref{prop:ss-level3-three-witness-reduction}
restricts to a nonzero equation on $\Lambda$.
Since its zero locus is a hypersurface,
the preceding divisorial-support calculation gives the asserted
containment.  Equivalently, any additional component of the
discriminant would have a divisorial singularity incidence, which
the chart calculation has exhausted.

Suppose a linear form
$\alpha F+\beta T+\gamma U+\delta v+\epsilon h$ vanished on
$\Delta_{\rm mov}$.  Substitute the displayed moving-node formulas.
Since $T,h$ vary freely, its coefficient of $h$ is identically zero:
\[
                  \gamma Z^2+\delta(z^2Z+w+d)+\epsilon=0.
\]
On the second elliptic curve the three functions
$z^2Z+w+d$, $Z^2$, and $1$ have respective pole orders $6,4,0$
at $O$.  Their leading terms cannot cancel, so
$\delta=\gamma=\epsilon=0$.
The coefficient of $T$ now gives
$\alpha\ell/A_1+\beta=0$ on a dense open of the product.
Fix a second point with $\ell\ne0$ and vary the first.
The function
$A_1=y^2+b+ax^2$ has pole order six and is nonconstant.
It follows that $\alpha=\beta=0$.
Thus no nonzero linear form vanishes on $\Delta_{\rm mov}$.

Finally a projective hyperplane contained in the singular-source
locus would, by irreducibility and the displayed support, equal one
of the four forced hyperplanes or equal $\Delta_{\rm mov}$.
The latter is impossible by the span assertion.  The complement
of the singular-source locus in every other hyperplane is therefore
a nonempty open.
\end{proof}

\begin{remark}[Exact symbolic checks]
\label{rem:DE-A4-third-value-checks}
The retained program
\path{agents/obstructions_crystals/code/contact_third_value_checks.py}
checks the universal moving-node equation, both first derivatives,
the nonzero mixed-Hessian formula, the coefficient of $h$ in an
arbitrary linear relation, the alternate formula for $A_1$,
$\det S=FTU$, and $l_{-g}=F+T$.
Its author run passed in $0.12054040400107624$ seconds; the complete
output is retained alongside the program.
Root and monodromy independently audited the entire proof and program,
with author-output-preserving reruns in $0.15617483800451737$ and
$0.1453385850036284$ seconds, respectively.
The pole-order independence and geometric incidence arguments are
proved in the text and are not inferred from the computation.
\end{remark}

\begin{corollary}[Every third-value hyperplane has smooth sharp sources]
\label{cor:DE-A4-third-value-smooth}
\status{Proved; root and monodromy full independent proof/code audits passed.}
Let $K\subset E^2(k)$ be a finite odd subgroup with both coordinate
projections injective, fix $0\ne g\in K$, and construct $\Lambda$
using $g_1,g_2$ as above.
For every $k\in K\setminus\{0,g\}$, the hyperplane
\[
                         H_k=\{s(k)=0\}\subset\Lambda
\]
has a nonempty open of globally smooth source curves with
$\det S\ne0$ and $Q_H(S)\ne0$.
Every such curve contains the three distinct points $P,Q,k$.
The assertion concerns smoothness and the source Hasse factors;
it does not claim that the three tangent directions have the
required tangent-triple type or that other norm collisions are absent.
\end{corollary}
\begin{proof}
In the contact coordinates the value functional at an affine point is
\[
      l_k=P_1(k)F+\ell(k)T+X(k)Z(k)U
                     +X(k)v+X(k)\ell(k)h,
      \qquad P_1=Y+a^2X.
\]
If $k=-g$, this is $F+T$.
Otherwise injectivity of both coordinate projections gives
$X(k)\ne0$ and $Z(k)\ne0$: equality of either corresponding
$x$-coordinate would force $k=g$ or $k=-g$.
Thus both the $U$ and $v$ coefficients of $l_k$ are nonzero.
In either case $H_k$ is distinct from the four forced hyperplanes
in Lemma~\ref{lem:DE-A4-contact-discriminant-span}, and from $U=0$.
That lemma gives a dense open of smooth sources in $H_k$.

The homogeneous determinant on the contact slice is
$\det S=FTU$, so its nonvanishing is another nonempty open on $H_k$.
The restricted source cubic $q=Q_H|_\Lambda$ is a nondegenerate
Frobenius cubic in five variables.  Its coefficient matrix has
determinant one, as in
Lemma~\ref{lem:ss-contact-cubic-evaluation}.
Polarization of this cubic excludes a projective plane in it,
and in particular excludes the hyperplane $H_k$.
Hence $q\ne0$ is also a nonempty open on $H_k$.
The irreducibility of $H_k\simeq\mathbf P^3$ lets these three opens
intersect.  The stated point containments are the defining contact
and third-value equations.
\end{proof}
